\subsection{片方零結合の射影条件からのOre認識}
\label{subsec:cubic-third-root-unit-onezero-projection}
\begin{proposition}[片方零結合の曲率のない固有直線商]
\label{prop:cubic-third-root-unit-onezero-pure-c}
$P=(1,0),Q=0,S=0$ の元22tensor atlasでは
\[
 B/(x,y,b)\simeq k[c].
\]
従ってH3候補正常三次 $G$ の純 $c^3$ 係数は零である。
$P=0,Q\ne0$ には $b,c$ を交換した主張が成り立つ。
\end{proposition}
\begin{proof}
重み $\omega$ の関係像を与えるmiddleの $c^2$ 行は零である。
他の重みの二次関係にも $c^2$ は現れない。
従って三生成元を零とした表示には $c$ の関係は残らない。
$Gc=\omega cG$ をこの商に写すと、純 $c^3$ 係数 $d$ に対し
$(1-\omega)d c^4=0$、従って $d=0$。
これはOre基底の存在を仮定しない議論である。
\end{proof}

\begin{theorem}[零even-cornerと可逆なmixed射影の全排除]
\label{thm:cubic-third-root-unit-onezero-projection}
$P=(1,0),Q=0,S=0,T=a_{21}=0$ とし、H3を与えると仮定する。
$U=\langle x,y\rangle$ とし、二次関係空間の重み $\omega^2$ 部分を
\[
 I_c\subset (U\otimes c)\oplus(c\otimes U)\oplus k b^2
\]
と書く。次の二条件が同時に成り立つ領域にはH3は存在しない：
\begin{enumerate}
\item $I_c$ から $U\otimes c$ と $c\otimes U$ への両射影が同型である。
\item even関係の $bc,cb$ への射影像を生成するベクトルの両係数が非零である。
\end{enumerate}
この領域の外、特に $T\ne0$ を本定理の結論に含めない。
\end{theorem}
\begin{proof}
元even blockは $U^2,bc,cb$ の分割で
$\left(\begin{smallmatrix}H&R\\R^t&0\end{smallmatrix}\right)$
という対称行列である。全rankは2なので $\operatorname{rank}R\le1$。
$R=0$ なら二つの独立な関係が純 $U^2$ に入り、
その二次元pencilはrank1 tensorを含むため、$B$ の整域性に反する。
従って $\operatorname{rank}R=1$。
全rank2から、even関係像は次の二本で書ける：
\[
 Q(U)=0,\qquad bc-qcb=h(U),\quad q\ne0,
\]
ここで最後の非零性には条件2を使う。
$Q$ は非零で、rank1なら積の零関係となるのでrank2である。
従って二元一関係代数 $T_0=k\langle x,y\rangle/(Q)$ は二次元AS-regularであり、
量子又はJordan標準形の語基底を持つ。
ここで $T_0\hookrightarrow B$ はまだ仮定していない。

重み $\omega$ の関係空間は $Ub\oplus bU$ に含まれる二次元空間である。
いずれかの射影に核があれば純 $ub$ 又は $bu$ の関係が生じ、整域性に反する。
従って両射影は同型であり、可逆線型写像 $f$ を用いて
$b u=f(u)b$ と書ける。
条件1により残りのmixed関係は
\[
 c u=g(u)c+d(u)b^2\quad(u\in U)
\]
と書け、$g$ は可逆である。$P\ne0$ なので $d\ne0$。

以上の六関係を使うと、全語は
$T_0$ の標準語 $\cdot b^i c^j$ で張られる。
実際、まず $c$ を右へ動かすと、補正項は $c$ の個数を下げる。
次いで $b$ を右へ動かし、最後に $T_0$ の関係で整理する。
この標準語の個数は各次数で $1/(1-t)^4$ の係数に一致する。
仮定したAS4 Hilbert級数も同じなので全語は独立であり、
$T_0$ と $T_0[b;f]$ は実際に埋め込まれる。
結合則を標準語係数で比較すると、$f$ は $T_0$ の自己同型となり、
$c$ の係数は $T_0[b;f]$ の自己同型 $\sigma$ と $\sigma$-導分 $\delta$ を与える。
$g$ と $q$ は可逆なので $\sigma$ も自己同型である。
従ってこれは実際の
\[
 B=T_0[b;f][c;\sigma,\delta],\quad
 \sigma|_{T_0}=g,\quad\sigma(b)=q^{-1}b,\quad
 \delta(u)=d(u)b^2,\quad\delta(b)=q^{-1}h
\]
という表示である（$h$ の符号は既存定理の記号へ吸収する）。
定理\ref{thm:cubic-third-root-unit-curved-Ore-empty}が適用でき、矛盾を得る。
これで、実Ore表示を前提とした既存排除へ、元tensorの明示射影条件から到達した。
\end{proof}

\begin{corollary}[片方零結合の零even-corner全排除]
\label{cor:cubic-third-root-unit-onezero-Tzero-empty}
$P=(1,0),Q=0,S=T=0$ の全層はH3を与えない。
従って重複1の未解決領域は、片方零結合かつ $T\ne0$ に限る。
\end{corollary}
\begin{proof}
前定理の二つの射影条件が、実は整域性と元middleのrank2から自動的に従うことを示す。
evenのcross二列を
\[
 R=(a_2,a_6,-a_{11},a_{15})^t,\qquad
 S=(a_3,a_7,-a_{13},a_{16})^t
\]
と書く。既に示した通り、cross rankは1である。

まず $R=0$ は不可能である。
元middle $M$ の行順を $xb,yb,bx,by,c^2$、
列順を $xc,yc,b^2,cx,cy$ とする。
$R=0$ なら最初の二行の最後の二列は零である。
$b$-mixedの $Ub$ 射影は整域性からrank2なので、この二行は $M$ の行空間の基底である。
従って全行の最後の二列が零となる。
第三・第四行から $a_8=a_9=a_{12}=a_{17}=0$ となり、
第二行全体が零となるので矛盾する。
$S=0$ の場合も、$bU$ 射影のrank2により第三・第四行を基底とし、
最初の二列の消滅から同じ四係数が零となり、第四行が零となる矛盾を得る。
従って $R,S$ は共に非零であり、$R=\rho S$、$\rho\ne0$ と書ける。
これは前定理のeven射影の両係数非零条件を与える。

純 $U$ 二次関係は $S$ で生成され、その係数行列は
\[
 Q=\begin{pmatrix}a_3&a_7\\-a_{13}&a_{16}\end{pmatrix}.
\]
整域性から $\det Q\ne0$。
元middleの二つのminorを直接取ると
\[
 \det M_{(bx,by),(xc,yc)}=(\omega+1)^2\det Q,\qquad
 \det M_{(xb,yb),(cx,cy)}=\rho^2\det Q.
\]
従って $c$-mixedの両射影も同型である。
前定理が適用でき、全 $T=0$ 層が排除される。
元22tensorの関係を使う二つのminor恒等式は、公開再現器でもexactに検算する。
\end{proof}
